\documentclass[UTF8,a4paper,12pt]{ctexart}

\usepackage{geometry}
\usepackage{booktabs}
\usepackage{amsmath}
\usepackage{amssymb}
\usepackage{listings}
\usepackage{xcolor}
\usepackage{hyperref}

\geometry{left=2.5cm,right=2.5cm,top=2.5cm,bottom=2.5cm}

\lstset{
    basicstyle=\ttfamily\small,
    keywordstyle=\color{blue},
    commentstyle=\color{gray},
    stringstyle=\color{teal},
    breaklines=true,
    frame=single,
    columns=fullflexible
}

\title{并行程序设计 Lab3：NTT 的 Pthread、CRT 与 OpenMP 优化实验报告}
\author{姓名：强徽东 {\hspace{3cm}} \quad 学号：2411764}
\date{\today}

\begin{document}

\maketitle

\section{实验背景与目标}

本实验选题为 NTT（Number Theoretic Transform，数论变换）下的多项式乘法优化。给定两个长度为 $n$ 的多项式 $A(x)$ 和 $B(x)$，目标是在模数 $P$ 下计算：
\[
C(x)=A(x)B(x)\pmod P
\]

实验按照由浅入深的顺序展开：

\begin{enumerate}
    \item 实现朴素卷积的 Pthread 并行版本，作为 baseline。
    \item 实现基础 NTT 的 Pthread 优化版本。
    \item 实现 CRT 合并的 NTT 版本，包括三模数和四模数版本。
    \item 实现 CRT 的 Pthread 优化版本。
    \item 进一步使用 OpenMP 实现自动调度优化。
    \item 在 OpenMP 版本基础上尝试使用 \texttt{\#pragma omp simd} 自动向量化。
\end{enumerate}

实验中特别注意平台约束：每个任务只申请 1 个 CPU（8 核），因此多线程版本均控制总工作线程数不超过 8，并尽量采用``最多 7 个子线程 + 主线程参与计算''的模式。

\section{朴素卷积的 Pthread 优化}

\subsection{串行朴素算法}

朴素多项式乘法的直接写法为：

\begin{lstlisting}[language=C++]
for (int i = 0; i < n; ++i) {
    for (int j = 0; j < n; ++j) {
        ab[i + j] = (ab[i + j] + a[i] * b[j]) % p;
    }
}
\end{lstlisting}

该算法复杂度为 $O(n^2)$。如果直接并行外层循环，不同线程可能同时写入同一个 \texttt{ab[i+j]}，需要加锁或使用原子操作，会引入较大的同步开销。

\subsection{按输出系数并行}

本实验采用按输出系数 $k$ 分块的方式：
\[
ab[k]=\sum_i a[i]b[k-i]\pmod P
\]

其中 $i$ 的范围满足：
\[
0\le i<n,\quad 0\le k-i<n
\]

因此：
\[
i\in [\max(0,k-n+1),\min(n-1,k)]
\]

每个线程负责一段连续的 $k$，不同线程写入不同的 \texttt{ab[k]}，因此不存在数据竞争，也不需要互斥锁。

\begin{lstlisting}[language=C++]
for (int k = task->l; k < task->r; ++k) {
    int begin = std::max(0, k - n + 1);
    int end = std::min(n - 1, k);

    u128 sum = 0;
    for (int i = begin; i <= end; ++i) {
        sum += (u128)a[i] * b[k - i];
        if ((i & 255) == 255) sum %= mod;
    }
    ab[k] = (u64)(sum % mod);
}
\end{lstlisting}

该版本作为 baseline，能够验证 Pthread 编程框架和任务划分方式，但由于复杂度仍为 $O(n^2)$，在 $n=131072$ 时运行时间较长。

\subsection{线程数量设计}

根据实验要求，线程总数控制在 8 以内。实际实现中：

\begin{lstlisting}[language=C++]
int worker_threads = threads - 1;
\end{lstlisting}

只创建 \texttt{threads-1} 个子线程，主线程也执行一段计算任务：

\begin{lstlisting}[language=C++]
plain_worker(&tasks[worker_threads]);
\end{lstlisting}

这样可以避免额外创建第 8 个子线程，同时让主线程承担计算任务。

\section{基础 NTT 的 Pthread 优化}

\subsection{NTT 基本流程}

NTT 将多项式卷积转化为点值乘法，整体流程如下：

\[
A \xrightarrow{\mathrm{NTT}} \hat{A},\quad
B \xrightarrow{\mathrm{NTT}} \hat{B}
\]

\[
\hat{C}[i]=\hat{A}[i]\hat{B}[i]\pmod P
\]

\[
C \xleftarrow{\mathrm{INTT}} \hat{C}
\]

复杂度由朴素卷积的 $O(n^2)$ 降为 $O(n\log n)$。

\subsection{蝶形循环并行}

基础 NTT 的核心循环为：

\begin{lstlisting}[language=C++]
for (int len = 2; len <= n; len <<= 1) {
    int half = len >> 1;
    for (int base = 0; base < n; base += len) {
        for (int k = 0; k < half; ++k) {
            // butterfly
        }
    }
}
\end{lstlisting}

同一层内不同蝶形块之间互不依赖，可以并行；但不同层之间存在依赖，因此每一层结束后需要同步。本实验中使用 \texttt{pthread\_barrier\_wait} 保证层间同步。

在前中期，蝶形块数量较多，按照 block 分配；在后期 block 数量少于线程数时，改为按 butterfly 数量划分，避免只有少数线程工作。

\section{CRT 实现任意模数 NTT}

\subsection{CRT 思想}

当目标模数 $P$ 过大，或者不是 NTT 友好模数时，可以选择若干个 NTT 友好小模数 $p_1,p_2,\ldots,p_m$。先分别在小模数下计算卷积：

\[
C_i(x)=A(x)B(x)\pmod {p_i}
\]

对于每个结果系数 $c[k]$，有：

\[
c[k]\equiv C_i[k]\pmod {p_i}
\]

使用中国剩余定理恢复 $c[k]\pmod M$，其中：

\[
M=\prod_{i=1}^{m}p_i
\]

最后再对目标模数 $P$ 取模：

\[
ans[k]=c[k]\bmod P
\]

CRT 正确恢复的前提是：

\[
n(P-1)^2 < M
\]

\subsection{三模数版本}

讲义中介绍的是三模数合并。本实验实现的三模数为：

\[
998244353,\quad 1004535809,\quad 469762049
\]

三者乘积约为：

\[
4.71\times 10^{26}
\]

该版本符合讲义模板，但对于特别大的目标模数，覆盖范围可能不足。

\subsection{四模数版本}

为了支持实验中出现的大模数，例如：

\[
P=263882790666241
\]

以及测试中出现的：

\[
P=1337006139375617
\]

进一步实现了四模数版本：

\[
998244353,\quad 1004535809,\quad 469762049,\quad 1224736769
\]

其中四个模数均为 NTT 友好模数，原根均可使用 3。四模数乘积约为：

\[
5.77\times 10^{35}
\]

能够覆盖 $n=131072$、$P=1337006139375617$ 时的最坏卷积系数上界。

\section{CRT 的 Pthread 优化}

CRT 版本天然具有粗粒度并行性。不同小模数下的 NTT 计算相互独立，因此可以并行执行。

四模数 Pthread 版本采用如下结构：

\[
4\text{ 个小模数任务}
=3\text{ 个子线程}+1\text{ 个主线程}
\]

即：

\begin{lstlisting}[language=C++]
const int worker_threads = CRT_CNT - 1;

for (int i = 0; i < worker_threads; ++i) {
    pthread_create(&tids[i], nullptr, mod_worker, &tasks[i]);
}

run_mod_worker(&tasks[worker_threads]);
\end{lstlisting}

小模数 NTT 全部完成后，再进行 CRT 合并。CRT 合并是按系数独立进行的，因此也可以按输出区间并行。

\section{OpenMP 自动调度优化}

\subsection{内部 NTT 并行版本}

在 OpenMP 版本中，首先实现了 NTT 内部并行：四个模数串行处理，但每个小模数内部的 NTT 蝶形循环使用 OpenMP 并行。

对于每层 NTT：

\begin{itemize}
    \item 当 block 数量较多时，使用 \texttt{schedule(static)}。
    \item 当 block 数量较少时，拆分 butterfly，避免线程空闲。
\end{itemize}

典型代码如下：

\begin{lstlisting}[language=C++]
#pragma omp parallel for schedule(static)
for (int block = 0; block < blocks; ++block) {
    // butterfly block
}
\end{lstlisting}

CRT 合并阶段使用：

\begin{lstlisting}[language=C++]
#pragma omp parallel for schedule(guided, 256)
for (int i = 0; i < result_len; ++i) {
    // CRT combine
}
\end{lstlisting}

\subsection{模数并行版本}

随后实现了另一种 OpenMP 优化方式：模数并行。外层四个小模数任务并行，每个小模数内部 NTT 保持串行，避免嵌套并行。

\begin{lstlisting}[language=C++]
int mod_threads = std::min(get_omp_threads(), CRT_CNT);

#pragma omp parallel for schedule(static) num_threads(mod_threads)
for (int i = 0; i < CRT_CNT; ++i) {
    multiply_mod_ntt(a, b, residues[i], n, CRT_MODS[i]);
}
\end{lstlisting}

这种方式最多使用 4 个线程，线程管理开销较小，但无法完全利用 8 核。对于 $n=131072$ 的大规模 NTT，内部 NTT 并行通常更容易充分利用核心；而模数并行作为对照实验，可以用于分析不同粒度并行策略的性能差异。

\section{OpenMP SIMD 自动向量化}

在模数并行 OpenMP 版本基础上，又实现了 OpenMP SIMD 自动向量化版本。

由于 NTT 内层原始写法中存在：

\[
w \leftarrow w\cdot wn
\]

这种循环依赖，不利于编译器向量化。因此先预计算旋转因子数组：

\begin{lstlisting}[language=C++]
std::vector<u64> roots(half);
roots[0] = 1;
for (int k = 1; k < half; ++k) {
    roots[k] = mul_mod_ntt(roots[k - 1], wn, mod);
}
\end{lstlisting}

之后在蝶形循环中使用：

\begin{lstlisting}[language=C++]
#pragma omp simd
for (int k = 0; k < half; ++k) {
    u64 x = f[base + k];
    u64 y = mul_mod_ntt(f[base + k + half], roots[k], mod);
    f[base + k] = add_mod(x, y, mod);
    f[base + k + half] = sub_mod(x, y, mod);
}
\end{lstlisting}

其中 \texttt{mul\_mod\_ntt} 使用 64 位乘法：

\begin{lstlisting}[language=C++]
static inline u64 mul_mod_ntt(u64 x, u64 y, u64 mod) {
    return x * y % mod;
}
\end{lstlisting}

由于 CRT 小模数均小于 $1.3\times 10^9$，乘积不超过 $1.6\times 10^{18}$，仍在 \texttt{uint64\_t} 范围内，因此该写法在小模数 NTT 内部是安全的。对于目标大模数下的通用乘法，仍应使用 \texttt{\_\_uint128\_t}。

\section{实验结果}

以下为部分已测试结果。时间单位以程序输出为准。

\begin{table}[h]
\centering
\caption{朴素 Pthread baseline 测试结果}
\begin{tabular}{cccc}
\toprule
样例 & $n$ & 模数 $P$ & 时间 / ms \\
\midrule
0 & 4 & 7340033 & 0.598961 \\
1 & 131072 & 7340033 & 7434.47 \\
\bottomrule
\end{tabular}
\end{table}

\begin{table}[h]
\centering
\caption{四模数 CRT 版本调试过程中的测试结果}
\begin{tabular}{cccc}
\toprule
样例 & $n$ & 模数 $P$ & 结果 \\
\midrule
0 & 4 & 7340033 & 正确 \\
1 & 131072 & 7340033 & 正确 \\
2 & 131072 & 104857601 & 正确 \\
3 & 131072 & 469762049 & 正确 \\
4 & 131072 & 1337006139375617 & 更换第四模数后正确 \\
\bottomrule
\end{tabular}
\end{table}

\section{问题与分析}

\subsection{三模数覆盖范围不足}

三模数 CRT 虽然符合讲义模板，但当目标模数较大时，可能无法覆盖真实卷积系数上界：

\[
n(P-1)^2
\]

因此在大模数测试中，需要使用四模数 CRT。实验中最初使用第四模数 $754974721$ 时，大模数样例出现错误；将第四模数替换为 $1224736769$ 后，四模数乘积增大，最终结果正确。

\subsection{线程粒度选择}

CRT 可以在模数层并行，也可以在 NTT 内部并行。模数并行粒度较粗，最多只有 4 个任务；NTT 内部并行粒度更细，能够更充分使用 8 核。因此对于大规模 $n=131072$，内部 NTT 并行通常更合适。

\subsection{SIMD 自动向量化限制}

\texttt{\#pragma omp simd} 只是向编译器提供向量化提示，并不保证一定生成 SIMD 指令。由于模乘中存在取模运算，编译器能否完全向量化取决于平台、编译器和优化参数。建议使用：

\begin{lstlisting}[language=bash]
g++ -std=c++17 -O3 -fopenmp -march=native main.cc -o main
\end{lstlisting}

\section{总结}

本实验围绕 NTT 多项式乘法实现了多层次优化。朴素 Pthread 版本验证了基本线程划分方法，但复杂度较高；基础 NTT 将复杂度降至 $O(n\log n)$；CRT 版本解决了大模数或非 NTT 友好模数的问题；Pthread 和 OpenMP 版本分别从手动线程控制和自动调度两个角度进行了并行优化；最后使用 OpenMP SIMD 尝试进一步利用向量化能力。

从实现过程看，CRT 模数数量和模数选择会直接影响大模数样例的正确性；线程数量需要结合平台核心数控制；OpenMP 的 \texttt{static}、\texttt{guided} 和 \texttt{simd} 可以分别用于规则循环、负载不均循环和内层向量化尝试。

\end{document}
